\section{معماری شبکه عصبی}
\label{sec:architecture}

تعریف معماری شبکه عصبی کانولوشنی شامل انتخاب نوع و تعداد لایه‌ها و تنظیم ابرپارامترهاست. این تنظیمات با هدف دستیابی به عملکرد رضایت‌بخش روی داده‌های اعتبارسنجی انجام می‌شوند؛ داده‌های آزمون باید برای ارزیابی نهایی مدل کنار گذاشته شوند. در این بخش، نوع و تعداد لایه‌های مورد استفاده در شبکه مشخص می‌شود. ساختار شبکه در این پروژه شامل لایه‌های زیر است:

\begin{itemize}
\item \textbf{لایه ورودی تصویر}\LTRfootnote{Image Input Layer}: داده‌های ولتاژ، ظرفیت دشارژ و دما به‌عنوان سه کانال یک تصویر ورودی در نظر گرفته می‌شوند و مقادیر آن‌ها در بازه $[0,1]$ نرمال‌سازی می‌شوند.
\item \textbf{لایه‌های کانولوشنی دوبعدی}\LTRfootnote{2-D Convolutional Layers}: هر یک از این لایه‌ها فیلترهای کانولوشنی را به‌صورت لغزشی روی تصویر ورودی اعمال می‌کنند. تعداد این لایه‌ها از طریق آزمون‌وخطا انتخاب شده است تا عملکرد مناسبی با زمان آموزش معقول فراهم شود.
\item \textbf{لایه‌های نرمال‌سازی لایه‌ای}\LTRfootnote{Layer Normalization Layers}: پس از هر لایه کانولوشنی، یک لایه نرمال‌سازی لایه‌ای قرار می‌گیرد. این لایه سرعت آموزش شبکه را افزایش می‌دهد و حساسیت آن را به مقداردهی اولیه کاهش می‌دهد.
\item \textbf{لایه‌های \lr{ReLU}}: پس از هر لایه نرمال‌سازی، یک تابع فعال‌سازی غیرخطی قرار می‌گیرد که عملگر آستانه‌گذاری را بر هر عنصر ورودی اعمال می‌کند.
\item \textbf{لایه‌های ادغام (پولینگ)}\LTRfootnote{Pooling Layers}: دو لایه نخست \lr{ReLU} با لایه‌های پولینگ همراه هستند که ابعاد نقشه ویژگی را کاهش می‌دهند. این لایه‌ها ضمن کاهش اطلاعات اضافی، تعداد پارامترهایی را که باید در لایه‌های بعدی یاد گرفته شوند نیز کم می‌کنند.
\item \textbf{لایه کاملاً متصل}\LTRfootnote{Fully Connected Layer}: در انتها، پس از لایه \lr{ReLU}، یک لایه کاملاً متصل قرار می‌گیرد که ویژگی‌های یادگرفته‌شده را در قالب یک خروجی واحد ترکیب می‌کند.
\end{itemize}

ساختار شبکه به‌ترتیب زیر تعریف شده است:






\begingroup
\renewcommand{\arraystretch}{3.2}
\begin{longtable}{|c|r|r|}
\caption{ساختار شبکه عصبی}\label{tab:layers}\\
\hline
\textbf{ردیف} & \textbf{لایه} & \textbf{پارامترها} \\ \hline
\endfirsthead
\hline
\textbf{ردیف} & \textbf{لایه} & \textbf{پارامترها} \\ \hline
\endhead
۱ & لایه ورودی تصویر & داده‌هایی با ابعاد $30\times30\times3$ و مقادیر نرمال‌شده در بازه ۰ تا ۱ \\ \hline
۲ & لایه کانولوشنی دوبعدی & ۸ فیلتر به اندازه $3\times3$ \\ \hline
۳ & لایه نرمال‌سازی & --- \\ \hline
۴ & لایه \lr{ReLU} & --- \\ \hline
۵ & لایه پولینگ بیشینه‌گیر & پنجره $2\times2$ \\ \hline
۶ & لایه کانولوشنی دوبعدی & ۱۶ فیلتر با اندازه $3\times3$ \\ \hline
۷ & لایه نرمال‌سازی & --- \\ \hline
۸ & لایه \lr{ReLU} & --- \\ \hline
۹ & لایه پولینگ میانگین‌گیر & پنجره $2\times2$ \\ \hline
۱۰ & لایه کانولوشنی دوبعدی & ۳۲ فیلتر با اندازه $3\times3$ \\ \hline
۱۱ & لایه نرمال‌سازی & --- \\ \hline
۱۲ & لایه \lr{ReLU} & --- \\ \hline
۱۳ & لایه کانولوشنی دوبعدی & ۳۲ فیلتر با اندازه $3\times3$ \\ \hline
۱۴ & لایه نرمال‌سازی & --- \\ \hline
۱۵ & لایه \lr{ReLU} & --- \\ \hline
۱۶ & لایه کانولوشنی دوبعدی & ۳۲ فیلتر به اندازه $3\times3$ \\ \hline
۱۷ & لایه نرمال‌سازی & --- \\ \hline
۱۸ & لایه \lr{ReLU} & --- \\ \hline
۱۹ & لایه کاملاً متصل & یک خروجی \\ \hline
\end{longtable}
\endgroup